%表紙
\newpage
\chapter{総則}
\begin{enumerate}[label=第\arabic*条　]
  \item この規則は、音楽ゲーム学園学園憲章（以下「憲章」という。）の趣旨に基づき、音楽ゲーム学園（以下「学園」という。）の組織及び運営に関し必要な事項を定めるものとする。
  \item この規則において、次の各号に掲げる用語の意義は、当該各号に定めるところによる。
      \begin{enumerate}[label=(\arabic*)　]
        \item 学生　第４章第３条の規定により学生としての資格を取得した者をいう。
        \item 教務主事　第２章第３条の規定により委嘱を受けた学生をいう。
        \item 学園長　第２章第４条の規定により学園長に充てられた教務主事をいう。
        \item 講師　第６章第２条の規定により認可を受けた学生をいう。
        \item 役職者　学園長、教務主事、講師その他の第２章第２条の規定により学園に置かれた役職にある者をいう。
        \item 講義　講師が学生に対して提供する学習活動をいう。
        \item 学期　学園における活動を区分する単位であって、第３章第２条に定める期間をいう。
        \item 正規活動期間　学園の講義その他の活動を通常の形態で実施する期間をいう。
        \item 補講期間　正規活動期間の終了後、当該学期における講義の補完的な実施のために設ける期間をいう。
        \item 募集期間　次の学期に向けて、新たに学生となろうとする者の募集を行う期間をいう。
      \end{enumerate}
  \item 憲章にいう構成員は、学生をいう。
  \item 憲章にいう運営は、教務主事をもって構成する。
  \item この規則の実施その他学園の運用に関し必要な事項を定めるため、ガイドラインを置くことができる。
    \begin{enumerate}[label=\arabic*　,start=2,leftmargin=*]
      \item ガイドラインは、憲章及びこの規則に反する内容を定めることができない。憲章又はこの規則に反するガイドラインは、その効力を有しない。
      \item ガイドラインは、この規則に定めのない義務を学生に課すことができない。
      \item ガイドラインの制定、改正及び廃止は、教務主事が行う。
      \item 教務主事は、学生の活動に関わるガイドラインを制定し、改正し、又は廃止したときは、速やかにこれを学生に公示しなければならない。
    \end{enumerate}
\end{enumerate}
